\submittedby{Manas Singh}

Let $\T_X$ denote the collection of all topologies on $X$, and let $\T_X/\mathrm{hom}$ denote the
set of homeomorphism classes of topologies on $X$, where $\T_1\sim\T_2$ if and only if
$(X,\T_1)$ and $(X,\T_2)$ are homeomorphic. Let $S_X$ be the group of all bijections
$f\colon X\to X$. Each $f\in S_X$ induces the bijection $\tilde f\colon\Pow(X)\to\Pow(X)$,
$\tilde f(A)=\{f(a)\mid a\in A\}$, and $\tilde f$ preserves arbitrary unions and intersections,
with $\tilde f(\emptyset)=\emptyset$ and $\tilde f(X)=X$.

\paragraph{The action is well-defined.} For $\T\in\T_X$, the collection
$f\cdot\T=\tilde f(\T)=\{f(U)\mid U\in\T\}$ contains $\emptyset$ and $X$ and is closed under
arbitrary unions and finite intersections, because $\tilde f$ preserves these operations. Hence
$f\cdot\T\in\T_X$. Moreover $\id\cdot\T=\T$ and $(g\circ f)\cdot\T=g\cdot(f\cdot\T)$, so this is
a group action of $S_X$ on $\T_X$.

\paragraph{Key observation.} For $f\in S_X$ and $\T_1,\T_2\in\T_X$,
\[ f\cdot\T_1=\T_2 \iff f\colon(X,\T_1)\to(X,\T_2) \text{ is a homeomorphism.} \]
Indeed, $f$ is a bijection, and $f\cdot\T_1=\T_2$ means precisely that
$U\in\T_1\iff f(U)\in\T_2$. The forward direction ($U\in\T_1\Rightarrow f(U)\in\T_2$) says
$f^{-1}$ is continuous, and the backward direction says $f$ is continuous. So $f$ is a
homeomorphism.

Conversely, every homeomorphism $f\colon(X,\T_1)\to(X,\T_2)$ is in particular a bijection
$X\to X$, i.e.\ an element of $S_X$, and it satisfies $f\cdot\T_1=\T_2$.

\paragraph{The bijection.} Define $\Phi\colon\T_X/S_X\to\T_X/\mathrm{hom}$ by
\[ \Phi(\operatorname{Orb}(\T))=[\T]_{\mathrm{hom}}. \]
\begin{itemize}
  \item \textbf{Well-defined and injective:} By the key observation,
  \begin{align*}
  \operatorname{Orb}(\T_1)=\operatorname{Orb}(\T_2)
  &\iff \exists f\in S_X:\ f\cdot\T_1=\T_2\\
  &\iff (X,\T_1)\cong(X,\T_2)
  \iff [\T_1]_{\mathrm{hom}}=[\T_2]_{\mathrm{hom}}.
  \end{align*}
  The implication ``$\Rightarrow$'' shows that $\Phi$ does not depend on the choice of
  representative of the orbit (well-defined), and ``$\Leftarrow$'' shows that $\Phi$ is
  injective.
  \item \textbf{Surjective:} Given any class $[\T]_{\mathrm{hom}}\in\T_X/\mathrm{hom}$, we have
  $\Phi(\operatorname{Orb}(\T))=[\T]_{\mathrm{hom}}$.
\end{itemize}
Therefore $\Phi$ is a natural bijection (it requires no choices) between the orbits of the
action of $S_X$ on $\T_X$ and $\T_X/\mathrm{hom}$.
